\section{Conclusion}
\label{sec:conclusion}

In this work, we investigated the complexity of model merging and proposed \textbf{Spectral Model Merging via Singular Value Slicing (SVS)}, a parameter-free, training-free, and closed-form merging operator. SVS performs a standard Singular Value Decomposition (SVD) on task-specific weight updates, slicing them to retain only the top $k$ principal singular components, and discarding the remaining spectral components as fine-tuning noise. We also analyzed \textbf{Barycentric Weight Normalization (BWN)}, a closed-form scale-preservation operator that matches the merged weights' Frobenius norm to the weighted barycenter of the individual experts, and discussed why this scale correction is mathematically neutralized in fully normalized networks.

Our empirical evaluation on CLIP-ViT-B/32 across four classification datasets (MNIST, FashionMNIST, CIFAR-10, SVHN) demonstrates the performance of SVS. SVS at rank $k=128$—retaining only $16.7\%$ of the rank space—matches or exceeds the performance of full-rank Task Arithmetic. Furthermore, in intermediate scaling regimes, SVS acts as an analytical regularizer, slightly exceeding standard Task Arithmetic. Operating entirely offline, SVS offers a zero-overhead merging solution with zero auxiliary parameters, zero test-time adaptation steps, and zero additional inference latency. Finally, we validated the scale-preservation behavior of BWN in un-normalized environments.

Future research directions include extending SVS to deeper transformer layers in multi-billion parameter Large Language Models (LLMs), conducting a comprehensive geometric sensitivity study of multi-dimensional tensor flattening choices (e.g., comparing input-channel versus output-channel grouping on singular value dispersion), and exploring hybrid merging paradigms that sequentially apply SVS as a continuous spectral filter followed by spatial coordinate-basis pruning to combine the complementary strengths of both methods. Ultimately, SVS demonstrates that model merging can be performed effectively through linear algebra principles without the need for auxiliary parameters or test-time optimization loops.
